% ===========================================
This section starts with a few words on Wilf's paper~\cite{Wilf1978AMM-circle}, including a transcript of the two problems that Wilf left open. Then, a little about each problem is said.
\subsection{Wilf's paper}
Wilf's concern was: 
to present an algorithm, which, given a numerical semigroup $S$ and a finite generating set for $S$, finds the conductor of $S$, decides whether a given integer is representable (in terms of the elements of the generating set), finds a representation of an element of $S$, and determines the number of omitted values of $S$.

These are problems that many researchers interested in combinatorial problems related to numerical semigroups are nowadays still concerned with. In a somewhat more modern language, one would say that Wilf was concerned with \emph{the Frobenius problem} (see~\cite{Ramirez-Alfonsin2005Book-Diophantine}), the membership problem, factorization problems (see~\cite{GeroldingerHalter-Koch2006}) and the problem of determining the genus.
These problems continue to be (are at the base of) active fields of research. 

The circle of lights algorithm explicitly given in~\cite{Wilf1978AMM-circle} determines both the conductor and the genus of a numerical semigroup provided that a finite generating set is at hand. Wilf also suggests a few changes to the circle of lights algorithm in order test membership and also to find a factorization. He also observed that to test membership, a suggestion of Brauer~\cite{Brauer1942AJM-problem} should be incorporated: it involves the use of the Apéry set (relative to the multiplicity).
Another observation that I would like to make is that the (space and time) complexity is explicitly given, which is a relevant contribution to the overall quality of Wilf's paper. It is extremely agreeable to read and this without doubt contributes to the success of the problems stated in it.

At the end of Wilf's article one finds the following two problems.
It should be understood that the positive integer $k$ represents the embedding dimension of some numerical semigroup. 

\begin{problem}[\cite{Wilf1978AMM-circle}]\label{prob:wilf-a-b}
	Wilf asked:
    \begin{itemize}
        \item[$(a)$] Is it true that for a fixed $k$ the fraction $\Omega/\chi$ of omitted values is at most $1-(1/k)$ with equality only for the generators $k,k+1,\ldots,2k-1$?
        \item[$(b)$] Let $f(n)$ be the number of semigroups whose conductor is $n$. What is the order of magnitude of $f(n)$ for $n\to\infty$?
    \end{itemize}
\end{problem}

The first problem consists in fact of two problems. They can be stated explicitly as follows:
\begin{problem}\label{prob:wilf-a1-a2} Wilf's problem (a) splits into two problems.
	\begin{itemize}
		\item[$(a.i)$] Is it true that for a fixed $k$ the fraction $\Omega/\chi$ of omitted values is at most $1-(1/k)$?
		\item[$(a.ii)$] Is it true that for a fixed $k$ the fraction $\Omega/\chi$ of omitted values is $1-(1/k)$ only for the generators $k,k+1,\ldots,2k-1$?
    \end{itemize}
\end{problem}

Problem ($a.i$) is nowadays known as \emph{Wilf's conjecture}.
Sylvester's result (which is mentioned in the first page of Wilf's paper) gives counter examples to Problem~$(a.ii)$. Apparently Wilf forgot about them. Nowadays there are other counter examples known, but a characterization of those semigroups for which the equality holds is an open problem (see Section~\ref{sec:wilf-a2}).

 %%%%%%%%%%
\subsection{Problem (\emph{a.i}): Wilf's conjecture}\label{sec:wilf-a1}
To a numerical semigroup~$S$ one can associate the following number denoted $\wilfoper(S)$ and called the \emph{Wilf number of~$S$}:
\begin{equation}\label{eq:wilf-number}
    \wilfoper(S) = \lvert \primitivesoper(S) \rvert\lvert \leftsoper(S) \rvert-\conductoroper(S).
\end{equation}

A numerical semigroup is said to be a \emph{Wilf semigroup} if and only if its Wilf number is nonnegative.
\emph{Wilf's conjecture} can be stated as follows:
\begin{conjecture}[Wilf, 1978]\label{conj:wilf}
  Every numerical semigroup is a Wilf semigroup.
\end{conjecture}
It is a simple exercise to verify that Conjecture~\ref{conj:wilf} is precisely Problem~\ref{prob:wilf-a1-a2}~($a.i$).
\medskip

There is another number that can be associated to a numerical semigroup, just as Wilf number is, and which has revealed great importance in recent research (as the reader will be able to confirm, in particular when reading Section~\ref{sec:large-mult}).
Let $S$ be a numerical semigroup and let $D= \primitivesoper(S)\cap \{\conductoroper,\dots,\conductoroper+\multiplicityoper-1\}$ be the set of non primitives in the threshold interval.
Eliahou~\cite{Eliahou2018JEMS-Wilfs} associated to $S$ the number $\eliahouoper(S)$ that appears in Equation~(\ref{eq:eliahou-number}) below and used the notation $\wilfoper_0(S)$ to represent it. I prefer the notation $\eliahouoper(S)$, and use the terminology \emph{Eliahou number of~$S$}:
\begin{equation}\label{eq:eliahou-number}
\eliahouoper(S) = \lvert \primitivesoper\cap \leftsoper\rvert\lvert \leftsoper \rvert - {\qoper} \lvert D\rvert + \qoper\multiplicityoper - \conductoroper.
\end{equation}

Eliahou~\cite[Pg.~2112]{Eliahou2018JEMS-Wilfs} observed that there are numerical semigroups with negative Eliahou number and stated the following problem which is still open.
\begin{problem}
	Give a characterization of the class of numerical semigroups whose Eliahou number is negative.
\end{problem}

%%%%%%%%%%
\subsection{Problem ($\emph{a.ii}$): another open problem}\label{sec:wilf-a2}
A very nice result of Sylvester~\cite{Sylvester1984ET-Mathematical} gives a formula for the Frobenius number of a numerical semigroup of embedding dimension $2$. A closed formula (of a certain type) for the Frobenius number of a numerical semigroup of higher embedding dimension is not at reach for semigroups of bigger embedding dimension (see~\cite{Curtis1990MS-formulas} or \cite[Cor. 2.2.2]{Ramirez-Alfonsin2005Book-Diophantine}). Sylvester's results can be written as follows (see~\cite{RosalesGarcia2009Book-Numerical}): if $S = \langle a, b\rangle$ is a numerical semigroup of embedding dimension $2$, then  $\Frobeniusoper(S) = ab-a-b$, and $\genusoper(S)=\conductoroper(S)/2$. From this, it is immediate that for a numerical semigroup $S$ of embedding dimension $2$, $\wilfoper(S)=0$.
The fact that numerical semigroups of the form $\langle \multiplicityoper, k\multiplicityoper+1,\ldots k\multiplicityoper+\multiplicityoper-1\rangle$ (which are of maximal embedding dimension and generated by some generalized arithmetic sequences) have Wilf number equal to $0$ is straightforward (see~\cite{FroebergGottliebHaeggkvist1987SF-numerical,MoscarielloSammartano2015MZ-conjecture}). 

Whether these are the only numerical semigroups for which Wilf number is $0$ is a slight modification of Problem~\ref{prob:wilf-a1-a2}($a.ii$) and is open. I rephrase the question stated by  Moscariello and Sammartano.

\begin{problem}[{\cite[Question~8]{MoscarielloSammartano2015MZ-conjecture}}]\label{prob:wilf-equality}
	\label{quest:W-equal-0}
	Let $S=\langle \multiplicityoper, g_2,\ldots,g_{\embeddingdimensionoper}\rangle$ be a numerical semigroup with multiplicity~$\multiplicityoper$ and embedding dimension~$\embeddingdimensionoper$. Is it true that if $\wilfoper(S)=0$, then $\embeddingdimensionoper(S)=2$ or $\embeddingdimensionoper(S)=\multiplicityoper(S)$ and there exists an integer $k>1$ such that $g_i=k\multiplicityoper+(i-1)$, for $i\in\{2,\ldots,\embeddingdimensionoper\}$? 
\end{problem}

Moscariello and Sammartano observed that in order to answer affirmatively this question it suffices to prove that for a semigroup with Wilf number equal to $0$, either its embedding dimension is $2$ or it has maximal embedding dimension.

They also observed that no numerical semigroup of genus up to $35$ provides a negative answer to the question.

Kaplan~\cite[Prop.~26]{Kaplan2012JPAA-Counting} has shown that Problem~\ref{prob:wilf-equality} has a positive answer in the case of numerical semigroups whose multiplicity is at least half of the conductor. The same holds for numerical semigroups of depth $3$ (see a remark by Sammartano in~\cite[Rem.~6.6]{Eliahou2018JEMS-Wilfs}), thus concluding that there are no counter examples among the semigroups satisfying $\conductoroper\le 3\multiplicityoper$.

%%%%%%%%%%%%
\subsection{Problem (b): counting numerical semigroups}

Wilf's Problem~\ref{prob:wilf-a-b}($b$) can be viewed as a problem about counting numerical semigroups by conductor. Backelin~\cite{Backelin1990MS-number} addressed this problem. A slight modification consists on counting by genus. Great attention has been given to this problem after Bras-Amorós~\cite{Bras-Amoros2008SF-Fibonacci} proposed some conjectures on the theme. Some of these conjectures were solved by Zhai~\cite{Zhai2012SF-Fibonacci}, while others remain open.
For an excellent survey (which in particular contains references for counting by conductor and has an outline of Zhai's proofs), see Kaplan~\cite{Kaplan2017AMM-Counting}.

Denote respectively by $N(g)$ and $t(g)$ the number of numerical semigroups of genus $g$ and the number of numerical semigroups of genus $g$ satisfying $\conductoroper(S) \le 3\multiplicityoper(S)$.

In the paper where he proved some of the conjectures of Bras-Amorós (one of them being that the sequence $\big(N(g)\big)$ behaves like the Fibonacci sequence), Zhai also proved that the proportion of numerical semigroups such that $\conductoroper(S) \le 3\multiplicityoper(S)$ tends to $1$ as $g$ tends to infinity, as conjectured by Zhao~\cite{Zhao2010SF-Constructing}.

\begin{proposition}[\cite{Zhai2012SF-Fibonacci}]\label{prop:Zhao-Zhai}
	With the notation introduced, the following holds:
	$$\lim_{g\to\infty} \frac{t(g)}{N(g)}=1.$$
\end{proposition}

I will leave here a question that can be stated in a similar way to Zhao's conjecture. It will be better appreciated when reading Section~\ref{sec:large-emb-dim} (and confronting with Section~\ref{sec:large-mult}).

Denote by $p(g)$ the number of numerical semigroups of genus $g$ satisfying $\embeddingdimensionoper(S) \ge \multiplicityoper(S)/3$.
\begin{question}\label{quest:emb-dim-asymptotics}
	Does $\lim_{g\to\infty} \frac{p(g)}{N(g)}$ exist?
\end{question}
% ===========================================


